% Generated draft: requires mathematical and source review.
\documentclass[11pt]{article}
\usepackage[margin=25mm]{geometry}
\usepackage{fontspec}
\setmainfont{Latin Modern Roman}
\usepackage{amsmath,amssymb,mathtools,unicode-math}
\setmathfont{Latin Modern Math}
\usepackage{enumitem,needspace,xurl,hyperref}
\setlength{\parindent}{0pt}
\setlength{\parskip}{5pt}
\setlength{\emergencystretch}{3em}
\newcommand{\N}{\mathbb{N}}
\newcommand{\Z}{\mathbb{Z}}
\newcommand{\Q}{\mathbb{Q}}
\newcommand{\R}{\mathbb{R}}
\newcommand{\C}{\mathbb{C}}
\newcommand{\F}{\mathbb{F}}
\newcommand{\A}{\mathbb{A}}
\newcommand{\PP}{\mathbb{P}}
\newcommand{\E}{\mathbb{E}}
\newcommand{\eps}{\varepsilon}
\newcommand{\dd}{\,\mathrm{d}}
\begin{document}
\section*{UnsolvedMath ID: 701495}
\textbf{Source problem code:} catalogue entry
\subsection{Definitions and mathematical statement}
For a group \(X\), define its derived series by \(X^{(0)}=X\) and \(X^{(i+1)}=[X^{(i)},X^{(i)}]\) for every integer \(i\geq 0\). If this series reaches the identity subgroup, the derived length \(\operatorname{dl}(X)\) is the least integer \(d\geq 0\) such that \(X^{(d)}=1\).

A finite semidirect product \(K\rtimes L\), with nontrivial kernel \(K\) and nontrivial complement \(L\), is a Frobenius group with that kernel and complement when \(C_K(\ell)=1\) for every \(1\neq\ell\in L\).

Let \(G\) be a finite group admitting a subgroup \(FH=F\rtimes H\leq\operatorname{Aut}(G)\) that is a Frobenius group with kernel \(F\) and complement \(H\). Thus \(C_F(h)=1\) for every \(1\neq h\in H\). Assume also that the semidirect product \(G\rtimes F\), denoted \(GF\) in the source, is a Frobenius group with kernel \(G\) and complement \(F\); equivalently, \(C_G(f)=1\) for every \(1\neq f\in F\). Define
\[
C_G(H)=\{g\in G: h(g)=g\text{ for every }h\in H\}.
\]
Is the derived length of \(G\) bounded solely in terms of \(|H|\) and the derived length of \(C_G(H)\)? Explicitly, does there exist a function
\[
b:\mathbb Z_{\geq 1}\times\mathbb Z_{\geq 0}\longrightarrow\mathbb Z_{\geq 0}
\]
such that, for every action satisfying these hypotheses and every integer \(d\geq 0\),
\[
C_G(H)^{(d)}=1\quad\Longrightarrow\quad G^{(b(|H|,d))}=1?
\]
This formulation expresses the requested bound through terminating derived series, without adding a separate solubility hypothesis.
\subsection{Short English statement}
Under the two stated Frobenius action hypotheses, must the derived length of \(G\) have a uniform bound depending only on the size of \(H\) and the derived length of the subgroup fixed by every element of \(H\)?
\subsection{Sources}
\begin{itemize}
\item Local UnsolvedMath registry, supplied complete record with permanent ID \(701495\), exact source problem code \texttt{catalogue entry}, and title ``Kourovka Notebook Problem 18.66.'' Supplied upstream revision: \texttt{d371ee917ac9647fbb06cd6fc529deec6f3dda82}.
\item The supplied background identifies \emph{The Kourovka Notebook (Unsolved Problems in Group Theory)}, 18th issue (2014), problem 18.66, and the 21st issue edition (2026), \texttt{arXiv:1401.0300v48}. These bibliographic details are supplied attribution and have not been independently checked here.
\item Proposed by N. Yu. Makarenko, E. I. Khukhro, and P. Shumyatsky; proposed year 2014, as recorded in the supplied registry.
\item The registry labels this problem ``open'' and records an update timestamp of \texttt{2026-10-16T14:00:00Z}. This is supplied registry information, not an independently verified assessment of the literature.
\end{itemize}
\end{document}
